
%% ================================================================
%%  8.  CONCLUSION AND FUTURE WORK
%% ================================================================
\section{Conclusion}\label{sec:conclusion}

This paper presented an open benchmarking framework for inventory
management that enables fair, reproducible comparison of solution
methods across four paradigms: heuristic policies, mathematical
programming, reinforcement learning, and deep learning. Rather than
creating a new simulation from scratch, we built upon and extended the
OR-Gym framework~\citep{hubbs2020orgym}---modernizing it to the
current Gymnasium~API, adding composable non-stationary demand
patterns and endogenous goodwill dynamics, and implementing classical
OR methods as first-class agents.

Our evaluation of ten solution methods across 32~controlled scenarios
yields several actionable insights:

\begin{enumerate}
  \item \textbf{No single paradigm dominates universally.} Classical
        OR methods (MSSP, DLP) excel under stationary demand with
        known parameters ($<6\%$ optimality gap), while GNN-based RL
        excels under endogenous demand dynamics (25\% improvement over
        MSSP) and provides 100--300$\times$ faster inference.

  \item \textbf{Endogenous dynamics are a key differentiator.} Methods
        that assume exogenous demand suffer significant degradation
        (25--45\% optimality gap) in goodwill-enabled scenarios.
        Learning-based methods that can implicitly capture
        demand--service feedback loops offer a genuine advantage.

  \item \textbf{Hybrid approaches are promising.} Residual~RL
        consistently improves upon its heuristic base, particularly
        in settings where the heuristic's assumptions are violated,
        while maintaining safety and interpretability.

  \item \textbf{The speed--quality trade-off is dramatic.} RL agents
        operate at two orders of magnitude faster than MSSP,
        achieving 85--95\% of MSSP's profit in non-goodwill settings,
        making them the only viable option for real-time applications.
\end{enumerate}

\subsection{Future Work}

The extensibility of the benchmarking framework opens several
promising research directions:

\paragraph{Environment Extensions.}
\begin{itemize}
  \item \textbf{Cost sensitivity analysis.} Systematically varying
        holding costs, backlog penalties, and prices to identify when
        the relative ranking of methods changes.
  \item \textbf{Stochastic lead times.} Introducing random lead-time
        variability, which creates a more realistic and challenging
        setting for all methods.
  \item \textbf{Multi-product settings.} Extending to environments
        with multiple products sharing capacity and exhibiting
        demand substitution, using our existing multi-agent
        extension based on PettingZoo.
  \item \textbf{Alternative demand distributions.} Testing
        heavy-tailed distributions (negative binomial, log-normal)
        and spatially correlated demand patterns.
\end{itemize}

\paragraph{Agent Extensions.}
\begin{itemize}
  \item \textbf{Foundation models.} Pre-trained time-series forecasters
        (e.g., TimesFM, Chronos) could serve as demand estimation
        modules feeding into any decision-making agent.
  \item \textbf{LLM-based agents.} Large language models could be
        evaluated as inventory decision-makers through prompt-based
        interaction with the Gymnasium environment.
  \item \textbf{Multi-agent RL.} Using the PettingZoo multi-agent
        extension to model decentralized decision-making with
        competing or cooperating supply chain entities.
  \item \textbf{Imitation learning.} Using DAgger or behavior
        cloning from Oracle trajectories to accelerate RL training.
\end{itemize}

\paragraph{Benchmarking Infrastructure.}
\begin{itemize}
  \item \textbf{Community leaderboard.} An online platform where
        researchers can submit agent implementations and compare
        results on the standardized scenario matrix.
  \item \textbf{Expanded scenario library.} Additional scenario
        configurations contributed by the community, including
        industry-specific cost structures and demand patterns.
\end{itemize}

The complete environment, all ten agent implementations, the
evaluation pipeline, and the data presented in this paper are
available as open-source software at
\url{https://github.com/r2barati/phd-paper1-draft}. We invite the
research community to build upon this framework, extend the scenario
library, and contribute new solution methods---advancing towards the
standardized benchmarking infrastructure that inventory management
research needs.

%% ================================================================
%%  ACKNOWLEDGEMENTS
%% ================================================================
\section*{Acknowledgements}

The authors acknowledge the computational resources provided by
Toronto Metropolitan University. This work builds upon the open-source
OR-Gym framework by \citet{hubbs2020orgym}.

%% ================================================================
%%  REFERENCES
%% ================================================================
\bibliographystyle{elsarticle-harv}
\bibliography{references}
